\subsection{Sine and Cosine Before Coordinates}

Take a right triangle containing an acute angle $\alpha$. Let $h$ denote the
hypotenuse, let $a$ be the leg adjacent to $\alpha$, and let $b$ be the leg
opposite $\alpha$.

\begin{center}
\begin{tikzpicture}[scale=0.95]
    \coordinate (A) at (0,0);
    \coordinate (B) at (4.8,0);
    \coordinate (C) at (4.8,2.9);
    \coordinate (D) at (3.0,0);
    \coordinate (E) at (3.0,1.8125);

    % Large right triangle
    \draw[subset] (A) -- (B) -- (C) -- cycle;
    \draw (B) ++(-0.3,0) -- ++(0,0.3) -- ++(0.3,0);

    % Smaller similar triangle cut out by a parallel segment
    \draw[orange!85!black,line width=1pt] (D) -- (E);
    \draw (D) ++(-0.22,0) -- ++(0,0.22) -- ++(0.22,0);
    \draw[orange!85!black,->] (3.13,0.63) -- (3.13,1.03);
    \draw[orange!85!black,->] (4.93,1.05) -- (4.93,1.45);

    % Common angle
    \draw (0.85,0) arc[start angle=0,end angle=31.1,radius=0.85];
    \node at (1.08,0.28) {$\alpha$};

    % Side labels for the large triangle
    \node[below] at (3.9,0) {$a$};
    \node[right] at (4.8,2.05) {$b$};
    \node[above left] at (3.9,2.35) {$h$};

    % Side labels for the smaller triangle
    \node[below] at (1.5,0) {$a'$};
    \node[left,orange!85!black] at (3.0,0.92) {$b'$};
    \node[above left] at (1.55,0.98) {$h'$};

    \node[orange!85!black,align=left,anchor=west] at (5.35,1.45)
        {$DE\parallel BC$\\[-1pt]\small Thales' theorem};
\end{tikzpicture}

\[
\frac{a'}{a}=\frac{b'}{b}=\frac{h'}{h},
\qquad
\frac{a'}{h'}=\frac{a}{h},
\qquad
\frac{b'}{h'}=\frac{b}{h}.
\]
\end{center}

For an acute geometric angle $\alpha$, define
\[
\cos\alpha=\frac{a}{h},
\qquad
\sin\alpha=\frac{b}{h}.
\]
These ratios are independent of the chosen right triangle because all right
triangles containing the same acute angle are similar.

For angles occurring in an arbitrary Euclidean triangle, we also need the right
and obtuse cases. We set
\[
\cos(\text{right angle})=0.
\]
If $\alpha$ is obtuse and $\beta$ is its acute supplement, so that together they
form a straight angle, define
\[
\cos\alpha=-\cos\beta.
\]
Thus the geometric cosine is defined for every angle between $0$ and a straight
angle. For acute angles it is positive, for a right angle it is zero, and for
obtuse angles it is negative.

The essential role of Thales' theorem and similarity is that these values depend
only on the angle, not on the size of the triangle used to represent it.

For an acute angle, the Pythagorean theorem gives
\[
a^2+b^2=h^2.
\]
Dividing by $h^2$ yields
\[
\left(\frac{a}{h}\right)^2+
\left(\frac{b}{h}\right)^2=1,
\]
and hence
\[
\cos^2\alpha+\sin^2\alpha=1.
\]
Thus the fundamental trigonometric identity is a direct reformulation of the
Pythagorean theorem.